\subsection{Inversos modulares}

Para achar o inverso multiplicativo de $a$ módulo $m$, uma condição necessária pra ele existir é que $\gcd(a, m)=1$. E como já falei em \ref{prob:exteuc1}, basta usar o algoritmo estendido de euclides.

Se não quiser usar ele, dá pra usar exponenciação binária pra calcular $a^{\phi(m)-1}$. Em particular, se $m$ for primo dá pra fazer isso em $a^{m-2}$.

\prob{1}{misc-tn} Ache inversos módulo $p$ primo.

\textbf{Solução 1:}

\begin{lstlisting}[language=c++, breaklines=true]
// acha inversos de 1 a n mod p primo
// acha em O(n) tempo. Requer O(n) memo

int n = 10, p = 1000000007;
int inv[n + 1];
inv[1] = 1;
for (int i = 2; i <= n; i ++) inv[i] = 1LL * (p - p / i) * inv[p % i] % p;
\end{lstlisting}

\textbf{Solução 2:}

\begin{lstlisting}[language=c++, breaklines=true]
// acha inversos de a[1], a[2], ..., a[n] mod p primo
// acha em O(n + log p) tempo. Requer O(n) memo

#include "bits/stdc++.h"
using namespace std;

const long long P = 1000000007;
long long qpow(long long a, long long b) {
	long long ans = 1;
	while (b) {
		if (b & 1) ans = ans * a % P;
		a = a * a % P;
		b >>= 1;
	}
	return ans;
}

long long n, a[1000010], pre[1000010], suf[1000010], pr = 1; // prefix, suffix, product
int main() {
	cin >> n; pre[0] = suf[n + 1] = 1;
	for (int i = 1; i <= n; i ++) cin >> a[i];
	for (int i = 1; i <= n; i ++) {
		pre[i] = pre[i - 1] * a[i] % P;
		suf[n + 1 - i] = suf[n + 2 - i] * a[n + 1 - i] % P;
		pr = pr * a[i] % P;
	}
	pr = qpow(pr, P - 2);
	for (int i = 1; i <= n; i ++) {
		cout << (pre[i - 1] * suf[i + 1] % P) * pr % P << endl;
	}
}
\end{lstlisting}
